\section{Introduction}

This paper considers prototype-based meta-learning, where models are trained to swiftly adapt to current tasks and perform classification through metric-based comparisons between examples and newly introduced variables - the prototypes of the classes. This approach, rooted in works by Reed \cite{reed1972pattern} and further developed by Snell \etal \cite{snell2017prototypical}, generalizes deep learning models to scenarios where labeled data is scarce. The fundamental idea is to compute the distances between queried examples and class prototypes and perform classification based on these distances to predict classes for queried examples. Derived from the prototypical network~\cite{snell2017prototypical}, several prototype-based methods have demonstrated their effectiveness in few-shot learning~\cite{allen2019infinite, ding2021prototypical, du2021hierarchical, zhen2020memory}. However, as prototypes are estimated from a limited number of sampled examples, they may not accurately capture the overall distribution~\cite{yang2021free, zhen2020memory}. Such biased distributions could lead to biased prototypes and subsequent classification errors, suggesting that the current prototype modeling approach might lack the ability to represent universal class-level information effectively.  Rather than relying on points in the prototype embedding space by using a simple average function, we propose to generate the distribution of prototypes for each task.



To overcome the challenges of estimating deterministic prototypes from limited examples in few-shot learning, we propose ProtoDiff, an innovative framework that leverages a task-guided diffusion model during the meta-training phase. ProtoDiff introduces a three-step process to address these limitations comprehensively. Firstly, we optimize a set of prototypes to achieve per-task prototype overfitting, ensuring accurate retrieval of overfitted prototypes specific to individual tasks. These overfitted prototypes serve as the \textit{ground truth} representations for each task. Secondly, a task-guided diffusion process is implemented within the prototype space, enabling meta-learning of a generative process that smoothly transitions prototypes from their vanilla form to the overfitted state. ProtoDiff generates task-specific prototypes during the meta-test stage by conditioning random noise on the limited samples available for the given new task. Finally, we propose residual prototype learning, which significantly accelerates training and further enhances the performance of ProtoDiff  by leveraging the sparsity of residual prototypes. The computational graph of ProtoDiff, illustrated in Figure~\ref{fig:graph}, showcases the sequential steps involved in the forward diffusion process on the prototype and the generative diffusion process. The resulting generated prototype, denoted as $\mathbf{z}_0$, is utilized for query prediction. By incorporating probabilistic and task-guided prototypes, ProtoDiff balances adaptability and informativeness, positioning it as a promising approach for augmenting few-shot learning in prototype-based meta-learning models.


\begin{figure*} [t!]
\centering
\includegraphics[width=0.98\linewidth]{graph/graph1.pdf}
\caption{\textbf{Computational graph of ProtoDiff.} In the initial stage, the computation of the vanilla prototype $\tilde{\mathbf{z}}$ is performed based on the support set $\mathcal{S}^{i}$. Subsequently, through the process of diffusion sampling, the diffused prototype $\mathbf{z}_{t-1}$ is derived from the combination of $\mathbf{z}_{t}$ and $\tilde{\mathbf{z}}$. Finally, the prediction $\mathbf{y}$ for the query set $\mathcal{Q}^{i}$ is generated by utilizing the diffused prototype $\mathbf{z}_0$ in conjunction with the query set $\mathbf{x}$. The diffusion forward and sampling processes are indicated by dashed and solid arrows within the gray rectangular area. }
	\label{fig:graph}
 \vspace{-2mm}
\end{figure*}


To validate the effectiveness of ProtoDiff, we conduct comprehensive experiments on three distinct few-shot learning scenarios: within-domain, cross-domain, and few-task few-shot learning. Our findings reveal that ProtoDiff significantly outperforms state-of-the-art prototype-based meta-learning models, underscoring the potential of task-guided diffusion to boost few-shot learning performance.
%
Furthermore, we provide a detailed analysis of the diffusion mechanism employed in ProtoDiff, showcasing its ability to capture the underlying data structure better and improve generalization. This thorough investigation highlights the strengths of our approach, demonstrating its potential to offer a more effective solution for few-shot learning tasks in various applications and settings.